\section{Experimental Evaluation}

\subsection{Experimental Setup}

The system is evaluated using a test set that covers heterogeneous enterprise-style queries:
\begin{itemize}
    \item Structured data queries handled by the tool path
    \item Document-grounded queries handled by the retrieval path
    \item General reasoning queries handled by direct generation
    \item Multi-tenant requests executed under isolated tenant contexts
\end{itemize}

Two configurations are compared:
\begin{itemize}
    \item Fixed single-path baselines, where queries are forced through one dominant processing pipeline
    \item The proposed adaptive system, which combines route selection, shared-model serving, and tenant-aware execution
\end{itemize}

\subsection{Evaluation Metrics}

The following metrics are used:
\begin{itemize}
    \item Latency (ms): time taken to produce a final response
    \item Memory usage (MB): RAM or GPU consumption during serving
    \item Response quality: correctness and usefulness of generated answers
    \item Path suitability: whether the selected execution path matches the query type
\end{itemize}

\subsection{Results and Analysis}

The results show that adaptive multi-path inference reduces latency by avoiding unnecessary retrieval or generation steps for lightweight queries. The shared-model design also lowers memory usage compared with configurations that replicate full models or overuse retrieval-heavy execution.

Structured queries benefit the most from routing because they can be resolved through tools without incurring the full cost of long-context LLM inference. Retrieval-heavy queries still benefit from grounding when needed, while general queries avoid unnecessary external lookup.

\subsection{Discussion}

Overall, the evaluation supports the central systems claim of the paper: adaptive route selection, shared-model serving, and lightweight personalization work together to improve efficiency while preserving multi-tenant practicality.
